\section{Avaliação e seleção de alinhamentos}
\label{sec:avaliacao-alinhamentos}

A previsão do resultado de uma equipa não fornece necessariamente uma
estimativa do desempenho de cada combinação de cinco jogadores.
Para apoiar a seleção de alinhamentos, é necessário distinguir duas
tarefas: estimar o valor das combinações e escolher uma combinação
admissível segundo os objetivos e as restrições. Os trabalhos de
Grassetti et al.~\cite{grassetti2021} e Kolias et al.~\cite{kolias2022}
contribuem sobretudo para a avaliação, enquanto Rodríguez Aguilar et
al.~\cite{rodriguez2020} exploram a otimização da seleção.

Grassetti et al.~\cite{grassetti2021} utilizam dados \emph{play-by-play}
de 240 partidas da EuroLeague 2018/2019 para desenvolver uma extensão
do \emph{Regularized Adjusted Plus-Minus}. A abordagem emprega
regularização Bayesiana e distribuições posteriores para analisar
contribuições individuais e coletivas e ordenar jogadores e
alinhamentos. Este enquadramento permite considerar o desempenho
conjunto, em vez de selecionar simplesmente os cinco jogadores com
melhores estatísticas individuais.

Kolias et al.~\cite{kolias2022} estudam a eficiência ofensiva e defensiva
de formações da Liga Grega, analisando 1\,131 fases de posse de bola
através de modelação estatística de Markov. A distinção entre as duas
dimensões de eficiência é relevante porque a preferência por um
alinhamento depende do objetivo: recuperar uma desvantagem pode exigir
uma escolha diferente da necessária para proteger uma vantagem.
Os dois estudos fornecem, assim, formas de caracterizar alternativas,
mas essa avaliação não equivale, por si só, a um algoritmo que escolha
a sequência de substituições futuras.

Rodríguez Aguilar et al.~\cite{rodriguez2020} apresentam um modelo
híbrido que combina otimização e avaliação de eficiência técnica na
seleção de alinhamentos da NBA. Segundo o resumo publicado, as
combinações selecionadas apresentaram maior eficiência técnica na
avaliação efetuada. Este resultado sustenta o interesse da programação
matemática para a composição da equipa; não estabelece automaticamente
um aumento da probabilidade de vitória nem resolve a gestão de minutos
ao longo de várias partidas.

A articulação destes trabalhos sugere uma interface entre os módulos
do projeto: a componente estatística estima o desempenho das alternativas
e a componente de otimização escolhe entre as alternativas admissíveis.
A qualidade da decisão depende de ambas. Uma solução ótima para uma
estimativa inadequada pode ter pouca utilidade desportiva; por isso,
a avaliação deve considerar a incerteza das estimativas e a
compatibilidade entre os dados disponíveis e as decisões representadas.

\section{Substituições e decisões sequenciais --- PLANTIA}
\label{plantia:estado-da-arte-de-plantia}

A seleção de um alinhamento num instante não resolve todo o problema
do treinador. As substituições alteram a composição da equipa e as
opções posteriores, pelo que importa representar a evolução dos estados
e comparar sequências de ações. A literatura analisada inclui modelos
descritivos de rotações, simulação de políticas e aprendizagem de
decisões a partir de dados históricos.

Bhat et al.~\cite{bhatSubstitutions} modelam substituições com dados
\emph{play-by-play} da fase regular da NBA de 2014/2015. Representam a
evolução dos alinhamentos por uma cadeia de Markov de tempo contínuo,
estimam taxas de transição e utilizam regressão para avaliar o
desempenho através do diferencial de pontos. Na simulação dos playoffs,
previram os vencedores de 12 das 15 séries. A dimensão reduzida desta
avaliação limita a generalização do resultado. Além disso, reproduzir
substituições observadas não demonstra que uma política alternativa
produziria melhores resultados.

Guo et al.~\cite{guo2023} analisam as rotações por uma perspetiva
complementar: redes de substituições construídas a partir de 574\,214
eventos de dez temporadas da NBA. O estudo identifica associações entre
a estrutura dessas redes e o desempenho competitivo. Enquanto Bhat et
al. representam transições entre alinhamentos, a análise de redes
caracteriza padrões de utilização e relações entre jogadores. Estes
indicadores podem apoiar a descrição das rotações, mas as associações
observadas não constituem uma política de decisão nem evidência causal
de que reproduzir determinado padrão melhora o desempenho.

Sandholtz e Bornn~\cite{sandholtz2020} avançam para a comparação de
políticas através de processos de decisão de Markov aplicados a
estratégias de lançamento. Combinam modelos Bayesianos hierárquicos,
probabilidades de transição dinâmicas e simulação, considerando a
evolução da posse e o tempo disponível para lançar. O interesse para
o projeto reside na representação explícita de estados, ações e
consequências, que permite avaliar escolhas alternativas. A adaptação
para alinhamentos, minutos e fadiga constituiu uma proposta do presente
trabalho, e não uma funcionalidade demonstrada nesse artigo.

Chen et al.~\cite{chen2022} exploram aprendizagem por reforço
\emph{offline} no sistema ReLiable, procurando aprender políticas
táticas a partir de registos históricos. Esta abordagem permite
investigar decisões sequenciais sem executar novas ações em jogos
reais. Constitui uma possível extensão do projeto, após a construção
e validação dos componentes fundamentais de simulação e avaliação.
A aprendizagem de uma política deve ser distinguida da demonstração
de que essa política é eficaz fora das condições em que foi avaliada.

Em conjunto, estes trabalhos oferecem elementos para descrever
rotações, estimar transições e comparar políticas. A escolha do método
depende da decisão pretendida e dos dados disponíveis: uma cadeia de
Markov pode representar a dinâmica observada, enquanto um processo de
decisão de Markov explicita as ações a comparar. No projeto, a passagem
de descrição para recomendação exigiu um ambiente comum, restrições
explícitas e estratégias de referência que permitiram avaliar o contributo
do planeamento.

\section{Síntese comparativa e oportunidades de investigação}
\label{plantia:anuxe1lise-cruxedtica-e-relauxe7uxe3o-com-o-projeto}

A Tabela~\ref{tab:comparacao-estado-arte} reúne os trabalhos por problema
e identifica os cuidados necessários à sua utilização no projeto.
As limitações sintetizadas dizem respeito à interpretação ou à
transferência dos contributos para o âmbito aqui proposto; não
representam necessariamente limitações declaradas pelos autores.

\begin{table*}[tp]
\centering
\small
\singlespacing
\caption{Comparação temática dos trabalhos e limites da sua aplicação ao projeto.}
\label{tab:comparacao-estado-arte}
\renewcommand{\arraystretch}{1.18}
\begin{tabular}{@{}
 >{\raggedright\arraybackslash}p{(\textwidth-6\tabcolsep)*\real{0.19}}
 >{\raggedright\arraybackslash}p{(\textwidth-6\tabcolsep)*\real{0.19}}
 >{\raggedright\arraybackslash}p{(\textwidth-6\tabcolsep)*\real{0.27}}
 >{\raggedright\arraybackslash}p{(\textwidth-6\tabcolsep)*\real{0.35}}@{}}
\toprule
Referência & Problema & Método ou contributo & Limite de aplicação ao projeto \\
\midrule
Lampis et al.~\cite{lampis2023} & Previsão pré-jogo & Comparação de ML e indicadores históricos & Superioridade dos algoritmos depende dos dados e da competição. \\
Horvat et al.~\cite{horvat2023} & Previsão na NBA & Construção de indicadores estatísticos & Prever resultados não define ações do treinador. \\
Ouyang et al.~\cite{ouyang2024} & Previsão explicável & XGBoost e SHAP & Explicações das previsões não estabelecem efeitos causais. \\
Kandhway~\cite{kandhway2025} & Previsão durante o jogo & ML supervisionado e SHAP & Informação disponível difere da previsão pré-jogo. \\
Li et al.~\cite{li2025} & Comparação da literatura & Revisão de 34 estudos & Dados e validações heterogéneos limitam comparações diretas. \\
\midrule
Grassetti et al.~\cite{grassetti2021} & Avaliação de alinhamentos & Plus-minus regularizado e modelos Bayesianos & Estimar desempenho não determina uma sequência de decisões. \\
Kolias et al.~\cite{kolias2022} & Eficiência de formações & Modelação estatística de Markov & Avaliação ofensiva e defensiva requer um planeador adicional. \\
Rodríguez Aguilar et al.~\cite{rodriguez2020} & Seleção de alinhamentos & Otimização e eficiência técnica & Eficiência técnica não equivale a vitórias ao longo de vários jogos. \\
\midrule
Bhat et al.~\cite{bhatSubstitutions} & Dinâmica de substituições & Markov em tempo contínuo e regressão & Avaliação em 15 séries; modelação observacional. \\
Guo et al.~\cite{guo2023} & Padrões de rotação & Análise de redes de substituições & Associações não demonstram vantagens causais de uma política. \\
Sandholtz e Bornn~\cite{sandholtz2020} & Decisões de lançamento & MDP com transições dinâmicas & Adaptação à gestão do plantel exige nova formulação. \\
Chen et al.~\cite{chen2022} & Políticas táticas & Aprendizagem por reforço offline & Transferência para gestão de minutos exige avaliação própria. \\
\bottomrule
\end{tabular}
\end{table*}

A primeira conclusão transversal é que a qualidade preditiva e a
qualidade das decisões são dimensões distintas. Um modelo pode prever
bem resultados agregados e não conter informação suficiente para
comparar alterações concretas de alinhamento. Do mesmo modo, uma
explicação estatística não quantifica necessariamente o efeito de uma
ação. A integração exige previsões compatíveis com as alternativas do
planeador e uma avaliação que distinga erros de previsão, cumprimento
de restrições e resultados das políticas simuladas.

A segunda conclusão diz respeito ao horizonte de decisão. Os trabalhos
analisados fornecem contributos para previsão, avaliação de alinhamentos,
substituições e políticas táticas, mas não demonstram diretamente uma
solução completa que combine seleção do cinco, distribuição de minutos,
fadiga e objetivos ao longo de várias jornadas. Esta observação delimita
uma oportunidade no conjunto de publicações analisado; não permite
concluir que inexistam outros sistemas integrados.

O contributo do trabalho consistiu em estudar essa articulação num protótipo
experimental, comparando decisões apoiadas por previsões com estratégias
fixas e heurísticas simples. O âmbito inicial privilegiou a previsão
pré-jogo e a simulação, mantendo a previsão dinâmica e a aprendizagem
por reforço como extensões possíveis. As questões de investigação e a
formulação encontram-se no Capítulo~\ref{chap:metodologia}; o protocolo
de comparação é apresentado no Capítulo~\ref{chap:experiencias}.
